\documentclass{scrartcl}
\usepackage{graphicx}  % required for inserting images

% for cyrillic and Bulgarian language
\usepackage[utf8]{inputenc}
\usepackage[T2A]{fontenc}
\usepackage[bulgarian]{babel}

% math packages
\usepackage{amsmath}
\usepackage{float}
\usepackage{amsfonts}
\usepackage{amssymb}
\usepackage{amsthm}
\usepackage{cancel}

\usepackage{ragged2e}  % adds FlushLeft
\usepackage{booktabs}
\usepackage{tabularx}

\usepackage{listings}
\usepackage{xcolor}

\definecolor{codegreen}{rgb}{0,0.6,0}
\definecolor{codegray}{rgb}{0.5,0.5,0.5}
\definecolor{codepurple}{rgb}{0.58,0,0.82}
\definecolor{backcolour}{rgb}{0.95,0.95,0.92}

\usepackage[margin=2cm]{geometry}

\lstdefinestyle{mystyle}{
    %caption=Example C++,
    label={lst:listing-cpp},
    language=C++,
    backgroundcolor=\color{backcolour},   
    commentstyle=\color{codegreen},
    keywordstyle=\color{magenta},
    numberstyle=\tiny\color{codegray},
    stringstyle=\color{codepurple},
    basicstyle=\ttfamily\footnotesize,
    breakatwhitespace=false,         
    breaklines=true,                 
    keepspaces=true,                 
    numbers=left,       
    numbersep=5pt,                  
    showspaces=false,                
    showstringspaces=false,
    showtabs=false,                  
    tabsize=2,
}
\lstset{style=mystyle}

\usepackage[colorlinks=true, linkcolor=blue, urlcolor=cyan]{hyperref}

\newtheorem{theorem}{Теорема}
\newtheorem{definition}{Дефиниция}

\title{Понятие за софтуерно инженерство. Софтуерен процес и модели на софтуерни процеси. Гъвкави методи за разработка на софтуер}
\author{Ивайло Андреев + Gemini 3 Flash}
\date{9 юни 2026}

\begin{document}
\maketitle

\tableofcontents

\newpage

\begin{FlushLeft}

\section{Понятие за софтуерното инженерство}

\subsection{Същност на софтуера и софтуерните системи}
За да дефинираме софтуерното инженерство, първо трябва да изясним същността на неговия обект на разработка — софтуера. Софтуерът не се ограничава единствено до изходния код или компютърните програми. Той представлява цялостна система, съставена от няколко ключови компонента:
\begin{itemize}
    \item \textbf{Компютърни програми:} Изпълнимият код (бинарни файлове, скриптове), стартиран от хардуера с цел постигане на определен резултат.
    \item \textbf{Структури от данни:} Логическите и физическите структури, които позволяват на програмите да манипулират, съхраняват и обработват информацията по ефикасен начин.
    \item \textbf{Документация:} Пълният набор от системни описания, изисквания, архитектурни дизайни, технически ръководства и потребителски инструкции, необходими за разбирането, използването и поддръжката на системата.
    \item \textbf{Интерфейси:} Потребителски (UI/UX) и програмни (API) интерфейси, които определят как системата си взаимодейства с външни потребители или с други софтуерни и хардуерни системи.
\end{itemize}

Софтуерът може да се разглежда в две основни направления: като самостоятелен \textbf{продукт}, предназначен за продажба или масово потребление, и като \textbf{средство за доставяне на услуги} (например SaaS модела), където фокусът е върху непрекъснатата стойност, предоставяна на потребителя.

\subsection{Характеристики на софтуера}
Софтуерът притежава специфични характеристики, които го отличават фундаментално от физическите продукти:
\begin{itemize}
    \item \textbf{Разработва се, не се произвежда:} Основните разходи и усилия в софтуерното инженерство са съсредоточени в етапа на проектиране и интелектуално създаване (R\&D). Самият процес на репликация (копиране на софтуера) е с практически нулева себестойност, за разлика от класическото производство.
    \item \textbf{Не се износва:} Софтуерът не страда от физическо износване или амортизация в традиционния смисъл. Вместо това той остарява морално поради промени в заобикалящата го среда (нов хардуер, операционни системи) или поради появата на нови изисквания. С времето, вследствие на многократни модификации, се проявява т.нар. „софтуерно остаряване“ (software deterioration), при което структурата му се усложнява.
    \item \textbf{Абстрактен продукт:} Софтуерът е неструктуриран физически и няма материална форма. Това го прави труден за визуализация, оценка на напредъка и концептуализиране.
    \item \textbf{Висока сложност:} Броят на възможните състояния в една средно голяма софтуерна система е на порядъци по-голям от този в коя да е механична система, което прави пълното тестване математически невъзможно.
    \item \textbf{Нужда от поддръжка и развитие:} Поради динамичната природа на бизнеса и технологиите, софтуерът изисква постоянни грижи за отстраняване на дефекти, адаптиране и добавяне на нови функционалности през целия му жизнен цикъл.
\end{itemize}

\subsection{Видове софтуер}
Съвременните софтуерни системи се класифицират в няколко основни категории в зависимост от тяхното предназначение и архитектура:
\begin{enumerate}
    \item \textbf{Системен софтуер:} Програми, написани да обслужват други програми (операционни системи, драйвери, компилатори).
    \item \textbf{Приложен софтуер:} Програми за крайни потребители, решаващи конкретни бизнес или лични задачи (текстообработващи програми, ERP системи, СУБД).
    \item \textbf{Научен/Инженерен софтуер:} Характеризира се с алгоритми за интензивна обработка на числа и симулации (CAD системи, астрономически изчисления).
    \item \textbf{Вграден софтуер (Embedded):} Намира се в паметта на хардуерни устройства и контролира техните функции (софтуер за автомобили, домакински уреди, медицинска апаратура).
    \item \textbf{Уеб приложения:} Системи, базирани на мрежова архитектура, които се достъпват през браузър и изпълняват логика на отдалечени сървъри.
    \item \textbf{Софтуер с изкуствен интелект (AI):} Използва неалгоритмични подходи, експертни системи, невронни мрежи и машинно обучение за решаване на комплексни проблеми.
    \item \textbf{Наследен софтуер (Legacy):} Стари, но критично важни за бизнеса системи, които са трудни за модификация, но не могат да бъдат лесно заменени.
\end{enumerate}

\subsection{Същност, цел и обхват на софтуерното инженерство}
\begin{definition}
\textbf{Софтуерно инженерство} е инженерна дисциплина, която се занимава с всички аспекти на производството на софтуер — от началните етапи на специфициране на системата до поддръжката и развитието ѝ след нейното въвеждане в експлоатация.
\end{definition}

Ключовите аспекти на дисциплината включват:
\begin{itemize}
    \item \textbf{Инженерна дисциплина:} Прилагане на систематичен, дисциплиниран и количествено измерим подход към разработката.
    \item \textbf{Методи, средства и процедури:} Използване на формални теории, инструменти (CASE средства) и утвърдени практики за управление на сложността.
    \item \textbf{Високо качество на софтуера:} Основен фокус върху надеждността, ефикасността, сигурността и използваемостта на крайния продукт.
    \item \textbf{Работа в рамките на срок и бюджет:} Инженерството винаги отчита икономическите и времевите ограничения на реалния свят.
\end{itemize}

\textbf{Основната цел} на софтуерното инженерство е създаването на \emph{качествен софтуер}, пълното \emph{удовлетворяване на изискванията на клиента} и поддържането на \emph{управляем и предвидим процес} на разработка.

\textbf{Обхватът} (lifecycle) обхваща следните фази: \textbf{Изисквания} (анализ и спецификация) $\rightarrow$ \textbf{Проектиране} (архитектура и детайлен дизайн) $\rightarrow$ \textbf{Реализация} (кодиране) $\rightarrow$ \textbf{Тестване} (верификация и валидация) $\rightarrow$ \textbf{Внедряване} (деплоймънт в реална среда) $\rightarrow$ \textbf{Поддръжка} (еволюция на системата).

---

\section{Софтуерен процес}

\subsection{Същност и елементи на софтуерния процес}
Софтуерният процес представлява \textbf{последователност от дейности}, събития и задачи, насочени към създаването, внедряването и еволюцията на даден софтуерен продукт. Всеки софтуерен процес се дефинира чрез пет основни стълба:
\begin{itemize}
    \item \textbf{Процес:} Самата структура и логика на преходите между различните фази.
    \item \textbf{Хора:} Ролите, отговорностите и екипната динамика на участниците в разработката.
    \item \textbf{Технологии:} Инструментите, програмните езици, фреймуорковете и хардуерната инфраструктура.
    \item \textbf{Ограничения и ресурси:} Финансови бюджети, времеви срокове (deadlines) и стандарти.
    \item \textbf{Междинни и крайни продукти:} Артефактите, които се генерират по време на процеса (модели, документи, софтуерни компоненти).
\end{itemize}

\subsection{Цели на софтуерния процес}
Добре дефинираният софтуерен процес преследва следните цели:
\begin{itemize}
    \item \textbf{Ефикасност:} Оптимално използване на човешките и технологичните ресурси.
    \item \textbf{Поддръжка:} Улесняване на бъдещите промени по системата чрез чиста структура.
    \item \textbf{Предсказуемост и повторяемост:} Възможност да се планират разходите и сроковете с висока точност на база предишен опит.
    \item \textbf{Качество:} Гарантиране, че дефектите се откриват възможно най-рано.
    \item \textbf{Усъвършенстване:} Постоянно анализиране и оптимизиране на самия начин на работа.
    \item \textbf{Проследяване (Traceability):} Възможност да се проследи всяка линия код до конкретно изискване на клиента.
\end{itemize}

\subsection{Основни дейности на процеса (Framework Activities)}
Всеки софтуерен процес, независимо от неговата конкретна специфика, включва пет основни дейности:
\begin{enumerate}
    \item \textbf{Комуникация (Communication):} Интензивен диалог със заинтересованите страни (stakeholders) с цел разбиране на целите на проекта и събиране на първоначалните изисквания.
    \item \textbf{Планиране (Planning):} Дефиниране на карта на пътя (roadmap), оценка на рисковете, разпределяне на ресурсите и определяне на графици.
    \item \textbf{Моделиране (Modeling):} Създаване на софтуерна архитектура, аналитични модели и дизайни (например чрез UML), които служат за основа на изграждането.
    \item \textbf{Конструиране (Construction):} Комбинация от кодиране (писане на изходния код) и тестване (unit и интеграционни тестове) за верификация на написаното.
    \item \textbf{Внедряване (Deployment):} Доставка на софтуерния инкремент/продукт до клиента, инсталация, конфигуриране и получаване на обратна връзка (feedback).
\end{enumerate}

\subsection{Допълнителни дейности (Umbrella Activities)}
Тези дейности се изпълняват паралелно през целия жизнен цикъл на проекта и служат за неговото осигуряване и управление:
\begin{itemize}
    \item \textbf{Управление на проекта:} Проследяване на напредъка спрямо плана.
    \item \textbf{Управление на риска:} Идентифициране и минимизиране на потенциални заплахи.
    \item \textbf{Осигуряване на качеството (QA):} Дейности за гарантиране на съответствието със стандартите.
    \item \textbf{Технически прегледи (Technical Reviews):} Проверки на кода и дизайна от колеги (peer reviews) с цел откриване на грешки.
    \item \textbf{Измерване (Measurement):} Събиране на метрики за процеса и продукта.
    \item \textbf{Управление на конфигурацията (SCM):} Контрол на версиите (например чрез Git) и промените по артефактите.
    \item \textbf{Повторно използване (Reusability):} Мениджмънт на компоненти за многократна употреба.
\end{itemize}

\subsection{Видове модели, езици за моделиране и шаблони}
Моделите на процеси могат да се разделят на две категории по предназначение:
\begin{itemize}
    \item \textbf{Описателни (Descriptive):} Описват как процесите се случват в реалността.
    \item \textbf{Предписателни (Prescriptive):} Предписват как трябва да се случват дейностите (действат като нормативни рамки).
\end{itemize}
Според перспективата на описание се дефинират: \textbf{Workflow Model} (последователност от действия), \textbf{Dataflow Model} (потоци от данни и трансформации) и \textbf{Role/Action Model} (кой какво прави).

За визуализация и специфициране се използва \textbf{UML} (Unified Modeling Language). Самите дейности в процесите често се описват чрез стандартизирани \textbf{шаблони}, съдържащи: \emph{Цел, Входни данни (Inputs), Изходни резултати (Outputs), Роли} и \emph{Ограничения}.

---

\section{Сравнителен анализ на описателни модели на софтуерен процес}

Описателните и предписателните модели дефинират конкретни структурни стратегии за управление и преминаване през жизнения цикъл на една система. Те се различават фундаментално по своята гъвкавост, отношение към управлението на риска, степен на документиране и въвлеченост на клиента.

\subsection{Модел на водопада (Waterfall Model)}
Моделът на водопада е първият утвърден и най-стар инженерен модел. Неговата основна философия се базира на \textbf{строга линейна последователност} – проектът протича като водопад отгоре надолу, като всяка фаза представлява затворен етап, който трябва да произведе фиксиран набор от документи и артефакти (Milestones).
\begin{itemize}
    \item \textbf{Специфика на протичане:} Едва след като фаза „Изисквания“ е напълно завършена, подписана от клиента и документирана, се преминава към „Проектиране“, след това към „Кодиране“ и т.н. Връщането към предходна фаза е изключително трудно и скъпо.
    \item \textbf{Предимства:} Висока степен на структурираност и яснота. Лесен за планиране и управление от мениджъри, тъй като крайните срокове и графиците за всяка фаза са дефинирани предварително.
    \item \textbf{Недостатъци (Критика):} Изключително неадаптивен към промени. В реалния бизнес клиентът рядко знае всичките си изисквания в самото начало. Най-големият риск е, че работещ софтуер се появява чак на финала на жизнения цикъл. Ако тогава се открие концептуална грешка от етапа на анализите, проектът може да претърпи провал.
    \item \textbf{Кога се използва:} При проекти с напълно ясни, дефинирани, стабилни и неизменни изисквания (например софтуер по държавни поръчки, банкови системи за транзакции или софтуер в авиацията).
\end{itemize}

\subsection{Прототипен модел (Prototyping Model)}
Създаден е с цел преодоляване на проблема с неясните потребителски изисквания във Водопада. Неговият фокус е върху \textbf{бързото изграждане на работещ макет (прототип)}, който се предоставя на клиента, за да може той да визуализира системата и да даде обратна връзка.
\begin{itemize}
    \item \textbf{Специфика на протичане:} Преминава през бърз анализ, проектиране на интерфейси, изграждане на прототипа и оценка от потребителя. Този цикъл се повтаря до изясняване на нуждите. Съществуват два основни подхода:
    \begin{enumerate}
        \item \emph{Еволюционен (Evolutionary):} Прототипът се пренаписва и рафинира стъпка по стъпка, докато се превърне в реалния краен продукт.
        \item \emph{Изхвърлен (Throw-away):} Прототипът служи единствено като средство за извличане на изисквания — често се използват мокове (mock обекти) вместо реална бизнес логика. След като изискванията се изяснят, той се изхвърля и софтуерът се изгражда на чисто с чиста архитектура.
    \end{enumerate}
    \item \textbf{Предимства:} Значително намалява риска от неразбирателство между разработчиците и клиента. Клиентът вижда как ще изглежда системата още в самото начало.
    \item \textbf{Недостатъци:} Бързото изграждане често води до компромиси с качеството на кода, сигурността и архитектурата. Често клиентите се подлъгват, че софтуерът е почти готов, и оказват натиск прототипът директно да бъде пуснат в експлоатация.
    \item \textbf{Кога се използва:} Когато клиентът има само обща представа за целите, но не знае детайлните бизнес правила, или при иновативни проекти с чисто нови потребителски интерфейси (UI/UX).
\end{itemize}

\subsection{Модел на бързата разработка (RAD - Rapid Application Development)}
RAD е модел, ориентиран към \textbf{изключително бързо и високоскоростно изграждане} на софтуерни системи (обикновено в рамките на 60 до 90 дни). Неговата философия е базирана на концепцията за паралелна работа и повторно използване на компоненти.
\begin{itemize}
    \item \textbf{Специфика на протичане:} Системата се разделя архитектурно на напълно независими модули. За всеки отделен модул се заделя самостоятелен екип. Всички екипи работят паралелно (едновременно), използвайки CASE инструменти за автоматично генериране на код и готови софтуерни компоненти. Накрая модулите се интегрират в обща система.
    \item \textbf{Предимства:} Изключително кратко време за излизане на пазара (Time-to-market). Спестява ресурси чрез повторна употреба на код.
    \item \textbf{Недостатъци:} Изисква огромни човешки ресурси (защото трябва да се поддържат множество паралелни екипи). Приложим е единствено за информационни системи, които са модулно делими. Ако архитектурата не позволява независимо разделение, моделът се проваля.
    \item \textbf{Кога се използва:} При големи бизнес информационни системи (например складови или административни софтуери), които трябва да се доставят в критично кратки срокове и има наличен голям брой опитни инженери.
\end{itemize}

\subsection{Спираловиден модел (Spiral Model)}
Предложен от Бари Боем (Barry Boehm), Спираловидният модел е еволюционен процес, който комбинира итеративната природа на прототипирането със системните контроли на модела на водопада. Неговата най-важна и отличителна характеристика е \textbf{акцентът върху управлението и анализа на риска}.
\begin{itemize}
    \item \textbf{Специфика на протичане:} Разработката се визуализира като спирала. Всяко завъртане (итерация) преминава през четири основни квадранта:
    \begin{enumerate}
        \item \emph{Определяне на целите} и алтернативните решения.
        \item \emph{Оценка на алтернативите и анализ на рисковете} (тук се правят прототипи и симулации за идентифициране на технически и бизнес заплахи).
        \item \emph{Разработка и верификация} на следващото ниво от продукта (същинско кодиране и тестване).
        \item \emph{Планиране} на следващата итерация и оценка от клиента.
    \end{enumerate}
    \item \textbf{Предимства:} Най-доброто управление на рисковете сред всички модели. Промените в изискванията се приемат естествено при всяко ново завъртане на спиралата. Продуктът се развива по контролиран начин.
    \item \textbf{Недостатъци:} Изключително сложен за управление и скъп модел. Успехът му зависи изцяло от експертизата на хората, които правят анализа на риска – ако даден критичен риск бъде пропуснат, моделът губи смисъла си.
    \item \textbf{Кога се използва:} При огромни, скъпи, комплексни и критични за сигурността системи (mission-critical софтуер, военни разработки, медицински хардуерно-софтуерни комплекси).
\end{itemize}

\subsection{Постъпков / Инкрементален модел (Incremental Model)}
Този модел комбинира елементи от линейния водопад с итеративната философия. Вместо софтуерът да се разработва наведнъж, той се \textbf{проектира, кодира и доставя на отделни парчета, наречени инкременти (постъпки)}.
\begin{itemize}
    \item \textbf{Специфика на протичане:} Първият инкремент задължително представлява ядрото на системата (Core product) – базовите функционалности, без които софтуерът не може. След като ядрото се внедри и тества от клиента, се разработват следващите инкременти (Инкремент 2, Инкремент 3), всеки от които добавя нови възможности и разширява системата.
    \item \textbf{Предимства:} Позволява на клиента да започне да използва основните функции на софтуера много рано. Рискът от пълен провал на проекта е нисък, тъй като ядрото се тества най-продължително.
    \item \textbf{Недостатъци:} Изисква перфектно софтуерно проектиране и архитектурно планиране в самото начало, за да се гарантира, че следващите инкременти ще могат лесно да се интегрират (вградят) към вече работещото ядро, без да го счупят.
    \item \textbf{Кога се използва:} Когато софтуерната система има ясна йерархия на функциите и бизнесът има нужда от бързо пускане на базова версия, докато пълната функционалност може да изчака във времето.
\end{itemize}

\subsection{Таблица за разширен сравнителен анализ}
Следващата таблица обобщава разликите между разгледаните описателни модели по ключови за изпита критерии:

\begin{table}[H]
\caption{Разширен сравнителен анализ на описателни модели на софтуерен процес}
\label{tab:detailed-process-models}
\small
\begin{tabularx}{\textwidth}{p{2.3cm} X X p{2.6cm}}
\toprule
\textbf{Модел} & \textbf{Управление на риска} & \textbf{Гъвкавост към промени} & \textbf{Участие на клиента} \\ 
\midrule
\textbf{Водопад} & Много слабо; рисковете се виждат чак в края при тестването. & Ниска; промени в късен етап изискват огромни разходи. & Основно в началото (изисквания) и в края (приемане). \\ \addlinespace
\textbf{Прототипен} & Добро по отношение на изискванията и интерфейса. & Висока в началния етап на дефиниране на макета. & Много силно по време на оценките на макетите. \\ \addlinespace
\textbf{RAD} & Средно; основният риск е интеграцията на модулите. & Средна; промени в обхвата на модула са трудни по време на спринта. & Умерено; във фазата на планиране и финално тестване. \\ \addlinespace
\textbf{Спираловиден} & Максимално високо; рискът се оценява на всяка итерация. & Много висока; всяко завъртане позволява адаптация. & Постоянно – в края на всяко завъртане дава обратна връзка. \\ \addlinespace
\textbf{Инкрементален} & Ниско до средно; предпазва ядрото на системата от провал. & Добра; новите изисквания се планират за следващ инкремент. & Регулярно – получава работещ софтуер след всяка постъпка. \\ 
\bottomrule
\end{tabularx}
\end{table}

---

\section{Гъвкави методи за разработка на софтуер (Agile)}

\subsection{Същност на Agile и Agile Manifesto}
Гъвкавите (Agile) методологии възникват като реакция срещу тромавите, силно документирани и тежки традиционни процеси (като Водопада). Те се базират на итеративен подход, бърза доставка на работещ софтуер, тясно сътрудничество с клиента и адаптивност към промените.

През 2001 г. се формулира \textbf{Agile Manifesto}, което дефинира 4 основни ценности и 12 принципа.
\textbf{Четирите ценности са:}
\begin{enumerate}
    \item \textbf{Хора и взаимодействия} пред процесите и инструментите.
    \item \textbf{Работещ софтуер} пред изчерпателната документация.
    \item \textbf{Сътрудничество с клиента} пред преговорите по договори.
    \item \textbf{Реакция на промяната} пред следването на предварително начертан план.
\end{enumerate}
Основните принципи включват непрекъсната доставка на стойност, техническо съвършенство, самоорганизиращи се екипи и постоянна адаптивност.

\subsection{Специфика на Scrum}
Scrum е най-популярната рамка (framework) за управление на гъвкави проекти. Тя дефинира три конкретни роли, три артефакта и набор от регулярни събития (практики):
\begin{itemize}
    \item \textbf{Роли:}
    \begin{itemize}
        \item \emph{Product Owner (Собственик на продукта):} Представлява гласа на клиента, управлява изискванията и приоритизира задачите.
        \item \emph{Scrum Master:} Служи като фасилитатор, премахва пречките (impediments) пред екипа и го защитава от външни смущения.
        \item \emph{Development Team (Екип за разработка):} Крос-функционален, самоорганизиращ се екип, който реално създава софтуера.
    \end{itemize}
    \item \textbf{Артефакти:}
    \begin{itemize}
        \item \emph{Product Backlog:} Динамичен, приоритизиран списък с всички нужни функционалности (епоси, потребителски истории, теми).
        \item \emph{Sprint Backlog:} Списък със задачи, избрани за изпълнение в текущия спринт.
        \item \emph{User Stories:} Изисквания, написани от гледна точка на крайния потребител върху карти.
    \end{itemize}
    \item \textbf{Практики:}
    \begin{itemize}
        \item \emph{Sprint:} Основен итеративен цикъл с фиксирана продължителност (обикновено от 2 до 4 седмици).
        \item \emph{Daily Scrum:} Кратка ежедневна 15-минутна среща, на която се споделя прогресът (какво е направено вчера, какво предстои днес и какви са проблемите).
        \item \emph{Работещ инкремент (Potentially Shippable Product Increment):} Резултатът от всеки спринт трябва да бъде работещ и тестван код.
    \end{itemize}
\end{itemize}

\subsection{Специфика на Extreme Programming (XP)}
Extreme Programming (XP) е софтуерна методология, която за разлика от Scrum (който е по-скоро управленски), се фокусира силно върху \textbf{инженерните и техническите практики} за писане на висококачествен код.
Ключовите инженерни практики на XP са:
\begin{itemize}
    \item \textbf{Pair Programming (Програмиране по двойки):} Двама програмисти работят заедно на едно работно място (един пише кода — Driver, другият го преглежда в реално време — Navigator).
    \item \textbf{Test-First Development (TDD):} Писане на автоматизирани unit тестове \emph{преди} имплементацията на самата функционалност.
    \item \textbf{Continuous Integration (Непрекъсната интеграция):} Често интегриране на готовия код в централното хранилище (няколко пъти на ден), последвано от автоматично пускане на всички тестове.
    \item \textbf{Refactoring (Рефакториране):} Постоянно подобряване на вътрешната структура на кода, без да се променя неговото външно поведение, с цел поддържане на простота.
    \item \textbf{Collective Ownership (Колективна собственост):} Всеки програмист от екипа има право да променя коя да е част от кодовата база по всяко време.
    \item \textbf{Simple Design (Прост дизайн):} Писане на код, който отговаря \emph{само} на текущите изисквания, без излишно презастраховане за бъдещето.
\end{itemize}

\subsection{Сравнителен анализ: Scrum срещу Extreme Programming}
Ето подробно съпоставяне на двете гъвкави методологии по ключови параметри на базата на предоставените източници:

\begin{table}[H]
\caption{Сравнение между Extreme Programming (XP) и Scrum}
\label{tab:xp-vs-scrum}
\small
\begin{tabularx}{\textwidth}{l X X}
\toprule
\textbf{Критерий} & \textbf{Extreme Programming (XP)} & \textbf{Scrum} \\ 
\midrule
\textbf{Фокус на метода} & Силно изразен фокус върху инженерните софтуерни практики. & Фокус върху управлението на проекта и организацията на екипа. \\ \addlinespace
\textbf{Итерации} & Кратки цикли; версиите могат да се билдват по няколко пъти на ден. Релийзи на всеки 2 седмици. & Фиксирани итерации (Спринтове), обикновено с продължителност 2–4 седмици. \\ \addlinespace
\textbf{Участие на клиента} & Изисква пълно работно ангажиране и постоянно присъствие на клиента в екипа. & Ролята на клиента се поема от Product Owner, който служи като посредник. \\ \addlinespace
\textbf{Управление на изискванията} & Инкрементно планиране чрез User Stories (карти), които се разбиват на задачи. & Product Backlog, съдържащ Epics, User Stories и Themes, подлежащи на приоритизация. \\ \addlinespace
\textbf{Инженерни практики} & Задължителни: Pair programming, TDD, рефакториране, колективна собственост. & Не предписва конкретни инженерни практики (често се комбинира с практиките на XP). \\ \addlinespace
\textbf{Промени в итерацията} & Допуска се замяна на задача в текущата итерация, ако не е започната и е със същия обем. & По време на Спринта изискванията са фиксирани, за да се защити екипът от промени. \\ \addlinespace
\textbf{Връзка с кода} & Поддържане на простота чрез постоянен рефакторинг и интеграция. & Продуктът се разделя на управляеми и разбираеми части (парчета софтуер). \\ 
\bottomrule
\end{tabularx}
\end{table}

\end{FlushLeft}
\end{document}